\documentclass[11pt,a4paper]{article}
\usepackage[T1]{fontenc}
\usepackage[utf8]{inputenc}
\usepackage{amsmath,amssymb,xcolor}
\usepackage[margin=25mm]{geometry}
\usepackage[colorlinks=true,urlcolor=blue]{hyperref}
\definecolor{nexorablue}{HTML}{2459B5}
\setlength{\parindent}{0pt}
\setlength{\parskip}{7pt}
\setlength{\emergencystretch}{2em}
\renewcommand{\familydefault}{\sfdefault}
\begin{document}

{\small\bfseries\color{nexorablue} NEXORA CPGE}\par
{\small MP / NEX-IND-04 / Énoncé}\par
{\LARGE\bfseries Puissance nécessaire pour tirer une barre sur des rails}\par
\vspace{0.25em}\hrule\vspace{0.6em}

Une barre de longueur $0{,}35\,\mathrm m$ se déplace à $2{,}5\,\mathrm{m\,s^{-1}}$ sur des rails dans un champ perpendiculaire de $0{,}32\,\mathrm T$. La résistance totale vaut $0{,}80\,\mathrm\Omega$. Calculer la force électromotrice, le courant et la force à exercer pour garder cette vitesse. Comparer les puissances électrique et mécanique.


\end{document}
